%----------------------------------------------------------------------------
\chapter{\osszefoglalas}
\label{ch:osszefoglalas}
%----------------------------------------------------------------------------

Jelen Tudományos Diákköri dolgozatban bemutattunk egy saját fejlesztésű, komplex videóalapú biomechanikai elemző és klasszifikációs szoftvercsővezetéket, amely saját rögzítésű kosárlabda büntetődobások kinematikáját vizsgálja. A megoldás alapját modern, marker nélküli pózbecslési technológiák (YOLO és MediaPipe) adják, amelyek jelentősen redukálják a mozgáselemzés eszköz- és költségigényét a klasszikus optikai mozgásrögzítő rendszerekhez képest.

A megvalósított szoftveres arhitektúra lefedi a kutatás teljes életciklusát a videók automatikus előfeldolgozásától és izolálásától kezdve az adatkinyerésig. Sikeresen integráltuk és összehasonlítottuk a TensorRT-vel GPU-n hardveresen felgyorsított YOLO-t, illetve a sokkal több referenciapontot (33 db) rögzítő, de CPU-ra szorítkozó MediaPipe-ot, amelynek teljesítményét kiterjedt aszinkron párhuzamosítással (parallelization) optimalizáltuk. Létrehoztunk egy valós idejű HUD vizualizációs réteget, valamint egyedi detektorokat fejlesztettünk az ízületi szögek robusztus meghatározásához.

Az előállított adathalmazra támaszkodva egy gépi tanulási klasszifikátor modellt építettünk fel, mely bizonyította, hogy a biomechanikai konzisztencia szorosan összefügg a dobások sikerességével. Bár az egyetlen monokuláris kamerára épülő elemzések korlátozottak a mikromozgások rögzítésénél, a modell megbízhatóan felismerte a hibás mozgásmintázatokat.

A kutatás eredményei rávilágítanak arra, hogy a marker nélküli pózbecslés – a GPU-s gyorsítások és a párhuzamosított architektúrák révén – elég érett technológia az élsportolói mozgáselemzés támogatásához. A jövőbeni továbbfejlesztések között szerepelhet a többkamerás rendszerek adatintegrációja a pontosság növelése érdekében, amellyel az edzők valós időben optimalizálhatnák a játékosok dobástechnikáját.
